\begin{pastexam}[label={px:ncu-107-dsa-18}]{中央資工所 107 學年度 第 18 題}
（複選；每題 5 分，答錯每小題倒扣 1 分，扣到零分為止）

Below is the Bellman-Ford Shortest Path (BFSP) algorithm, whose input is a positively weighted
digraph, and whose output is the distance (or accumulated weight) of the shortest path from the
source node to every node. Choose all items that all together make the algorithm a negative cycle
detection algorithm, whose input is a weighted digraph $G$ in which weights may be negative or
positive, and whose output is ``\,`yes' if graph $G$ has cycles of negative weights; otherwise,
`no'\,''.\\
(a)~replace ``$w[x][y]>0$'' with ``$w[x][y]>0$ or $w[x][y]<0$'' in the \textbf{Input} section\\
(b)~replace the \textbf{Output} with ``\,`yes' if graph $G$ has cycles of negative weights;
otherwise, `no'\,''\\
(c)~replace line 2 with\\
\hspace*{2em}2. \textbf{for} $i\leftarrow 1$ to $|V|$ \textbf{do}\\
(d)~replace line 6 with the following lines:\\
\hspace*{2em}6. \textbf{if} $d[u]>0$ for each $u\ne s$ \textbf{then}\\
\hspace*{2em}7. \qquad \textbf{return} ``yes''\\
\hspace*{2em}8. \textbf{else}\\
\hspace*{2em}9. \qquad \textbf{return} ``no''\\
(e)~replace line 6 with the following lines:\\
\hspace*{2em}6. \textbf{for} every edge $(x,y)$ in $E$ \textbf{do}\\
\hspace*{2em}7. \quad \textbf{if} $d[y]>d[x]+w[x][y]$ \textbf{then}\\
\hspace*{2em}8. \qquad\quad \textbf{return} ``yes''\\
\hspace*{2em}9. \textbf{return} ``no''
\begin{quote}
\textbf{Algorithm} BFSP\\
\textbf{Input:} a weighted digraph $G=(V,E)$, where the weight of each edge $(x,y)$ in $E$ is
represented by $w[x][y]$, $w[x][y]>0$, and a source node $s$, $s\in V$\\
\textbf{Output:} the distance of the shortest path from the source node to every node\\
1.\quad $d[s]\leftarrow 0$;\quad $d[u]\leftarrow\infty$ for each $u\ne s$\\
2.\quad \textbf{for} $i\leftarrow 1$ to $|V|$-1 \textbf{do}\\
3.\quad\quad \textbf{for} every edge $(x,y)$ in $E$ \textbf{do}\\
4.\quad\quad\quad \textbf{if}\ \ $d[y]>d[x]+w[x][y]$ \textbf{then}\\
5.\quad\quad\quad\quad $d[y]\leftarrow d[x]+w[x][y]$\\
6.\quad \textbf{return} $d$
\end{quote}
\end{pastexam}

\begin{pastsoln}
答案：(a)(b)(e)。

這就是資結書第~25.3~節 Bellman--Ford 的做法（虛擬碼~25.4）：鬆弛 $|V| - 1$ 輪之後，再把每條邊檢查
一次，還能鬆弛就表示有負環（資結書定理~25.7、第~25.3.4~節）。

\begin{itemize}
\item (a) 選。 新的問題要處理負權重，Input 的限制 $w[x][y] > 0$ 要放寬。 照字面它排除了權重 $0$，嚴格說
寫成「任意實數」比較好，但在五個選項裡這一項是必要的。
\item (b) 選。 輸出從距離改成 yes／no，Output 的描述要跟著改。
\item (c) 不選。 多鬆弛一輪不必要：沒有負環時，$|V| - 1$ 輪之後所有 $d$ 已經是最短距離（最短路徑最多
$|V| - 1$ 條邊），偵測負環靠的是 (e) 那一輪檢查，不是多跑一輪。 單獨改 (c) 而沒有檢查，演算法仍然只
回傳 $d$。 （加上 (c) 再配 (e) 結果也還是對的，只是多做一輪；題目要的是讓它變成負環偵測的那幾項。）
\item (d) 不選。 「所有 $d[u] > 0$ 就回答 yes」和負環沒有關係。 反例：只有一條邊 $s \to u$，權重 $1$，
沒有任何環；$d[u] = 1 > 0$，(d) 回答 yes，錯。 反過來，有負環時 $d$ 值會被拉成負的，(d) 反而回答 no。
\item (e) 選。 沒有負環時，$|V| - 1$ 輪後 $d[y] = \delta(s, y)$，三角不等式 $d[y] \le d[x] + w[x][y]$ 對
每條邊都成立，第 7 行不會成立，回答 no。 有從 $s$ 走得到的負環 $v_0, v_1, \ldots, v_k = v_0$ 時，若每條環
上的邊都滿足 $d[v_i] \le d[v_{i-1}] + w[v_{i-1}][v_i]$，沿環加起來得
$\sum d[v_i] \le \sum d[v_{i-1}] + \sum w$，兩邊的 $\sum d$ 相同（都是有限值），得 $0 \le \sum w$，和負環矛盾；
所以一定有一條邊能鬆弛，回答 yes。
\end{itemize}

補充：這個演算法只偵測得到\textbf{從 $s$ 走得到}的負環。 走不到的點 $d$ 一直是 $\infty$，那裡的負環檢查
不出來。 題目要「$G$ 有沒有負環」，嚴格的做法是加一個新起點連到每個點（權重 $0$），或一開始把所有
$d[u]$ 都設成 $0$。 這不在選項裡，照題意選 (a)(b)(e)。
\end{pastsoln}
